\section{Scope and limitations}
\label{sec:limitations}

The first limitation is semantic scope and external correspondence: the formal
B semantics covers the mathematical fragment \beer{} requires, and machines,
substitutions, refinement, and proof-obligation generation lie outside
its scope. For the fragment it does
cover, its agreement with the reference account of the
B-Book~\cite{abrial1996bbook} is argued informally, and necessarily so: that
account is written in prose, so the correspondence between it and any
formalization is a mathematical judgment made on paper. Every development of
this kind keeps that judgment in its trusted base. The SMT-LIB formalization is
likewise restricted to the fragment the encoder emits; formalizing the whole
standard would be a disproportionate effort for a guarantee the encoder never
exercises. The soundness theorem accordingly relates the formal source and
target models developed in this manuscript.

The second limitation is coverage: the verified \beer{} encoder leaves out, among
other constructs, cardinality, minimum, maximum, and substantial parts of the
sequence language. Unsupported inputs are rejected, which preserves soundness at
the cost of applicability. \textsc{Zero-Proof} \beer{} covers these constructs
and the rest of the extended fragment listed in \Cref{sec:beer-zero-proof}, and
is the encoder measured in \Cref{ch:evaluation}. It descends from the verified
encoder and follows the same encoding design; what falls outside the theorem is
the added fragment, together with the optimizations of the encoder and of the
emitted script that \Cref{sec:beer-zero-proof} lists.

\barrel{} supports a fragment of comparable extent, sufficient for the
developments studied here and open to extension: each missing construct is added
by the same short sequence of steps (\Cref{subsec:barrel-coverage}). Outside
this manuscript, the tool has drawn the attention of the industrial B community,
and projects to integrate it into existing toolchains are under way.